% !TeX program = xelatex
\documentclass[12pt]{article}
\usepackage[spanish, es-nodecimaldot]{babel}

\usepackage{fontspec}
\setmainfont{Times New Roman}

\usepackage[letterpaper,top=2.4cm,bottom=2.4cm,left=2.4cm,right=2.4cm]{geometry}

\usepackage{setspace}
% Renglon cerrado
\singlespacing
\usepackage{xurl}
\usepackage[hidelinks]{hyperref}

\usepackage{enumitem}
\usepackage{booktabs}
\usepackage{float}
\usepackage{tabularx}
\usepackage{array}

\usepackage{titlesec}

\setlength{\parindent}{0pt}

\setlength{\parskip}{6pt}

\titleformat{\section}
{\normalfont\bfseries\fontsize{12}{14}\selectfont}
{}{0pt}{}
\titlespacing*{\section}{0pt}{6pt}{0pt}

\author{Jenny Sofía Morales}

\begin{document}
	
	\begin{center}
		
		{\bfseries\fontsize{12}{14}\selectfont
			CRITERIOS PARA ELEGIR UNA ARQUITECTURA
		}
		
		\vspace{0.4cm}
		
		{\itshape J. S. Morales López}
		
		\vspace{0.2cm}
		
		{\itshape 7690-08-6790 Universidad Mariano Gálvez}
		
		{\itshape Seminario de tecnologías de información}
		
		{\itshape
			\href{mailto:jmorales115@miumg.edu.gt}
			{jmorales115@miumg.edu.gt}
		}
		
		{\itshape 
			\href{https://github.com/sofiaMorales3805/ForoAcademico-Seminario}
			{Repositorio en GitHub - Sofia Morales}
		}
	\end{center}
	
	\vspace{0.5cm}
	\section*{Resumen}
	
	El presente artículo analiza los principales criterios que deben considerarse
	para seleccionar la arquitectura de una aplicación. Esta decisión no depende
	únicamente de si el software será web, móvil o de escritorio, sino también de
	su complejidad, cantidad de usuarios, volumen de información, integraciones,
	disponibilidad requerida, presupuesto y experiencia del equipo de desarrollo.
	Se describen arquitecturas como la monolítica, en capas, de microservicios,
	orientada a eventos y sin servidor. Asimismo, se relacionan estas alternativas
	con aplicaciones pequeñas, sistemas empresariales, plataformas de gran escala,
	soluciones de Internet de las cosas y aplicaciones con demanda variable. El
	análisis demuestra que cada arquitectura presenta ventajas, limitaciones y
	costos técnicos que deben evaluarse antes de iniciar el desarrollo. Una
	arquitectura sencilla puede ser suficiente para un proyecto pequeño, mientras
	que un sistema complejo podría requerir componentes independientes y
	escalables. Se concluye que la arquitectura adecuada es aquella que satisface
	los requisitos actuales del proyecto, permite su mantenimiento y puede
	evolucionar sin introducir una complejidad innecesaria.
	
	\noindent
	\textbf{Palabras clave:} arquitectura de software, aplicaciones, microservicios, escalabilidad y requisitos.
	
	
	\section*{Desarrollo del tema}
	
	La arquitectura de software representa la organización general de una
	aplicación. Describe sus componentes principales, las responsabilidades de
	cada uno y la manera en que intercambian información. Entre estos componentes
	pueden encontrarse la interfaz de usuario, las reglas de negocio, las bases de
	datos, los servicios y las integraciones con sistemas externos.
	
	La selección de una arquitectura influye en el rendimiento, la seguridad, la
	mantenibilidad y la capacidad de crecimiento del sistema. Por esta razón, no
	debe seleccionarse solamente porque sea moderna o porque haya sido utilizada
	por una empresa reconocida. Según Kazman et al. (2023), las decisiones
	arquitectónicas deben relacionarse con las necesidades del sistema y con sus
	atributos de calidad. Esto significa que una arquitectura apropiada para una
	plataforma de comercio electrónico de gran tamaño puede resultar innecesaria
	para una aplicación académica utilizada por pocas personas.
	
	\section*{Criterios para elegir una arquitectura}
	
	Antes de seleccionar una arquitectura, se deben identificar los requisitos
	funcionales y no funcionales. Los requisitos funcionales establecen lo que
	debe realizar la aplicación, mientras que los no funcionales determinan
	condiciones como seguridad, rendimiento, disponibilidad, facilidad de
	mantenimiento y escalabilidad.
	
	Los principales criterios de selección son los siguientes:
	
	\begin{itemize}
		
		\item \textbf{Tamaño y complejidad:} se debe determinar cuántos módulos,
		procesos y reglas de negocio tendrá la aplicación.
		
		\item \textbf{Cantidad de usuarios:} permite estimar la carga de trabajo
		que deberá soportar el sistema.
		
		\item \textbf{Volumen de información:} ayuda a definir las necesidades de
		almacenamiento, procesamiento y consulta de datos.
		
		\item \textbf{Disponibilidad:} establece si la aplicación puede detenerse
		temporalmente o debe funcionar de manera continua.
		
		\item \textbf{Integraciones:} identifica la necesidad de comunicarse con
		servicios externos, dispositivos, API u otros sistemas empresariales.
		
		\item \textbf{Seguridad:} determina la sensibilidad de la información y
		los controles necesarios para protegerla.
		
		\item \textbf{Presupuesto y tiempo:} condiciona la infraestructura, las
		herramientas y la complejidad que puede asumir el proyecto.
		
		\item \textbf{Experiencia del equipo:} permite comprobar si los
		desarrolladores poseen los conocimientos necesarios para implementar y
		administrar la arquitectura.
		
		\item \textbf{Crecimiento esperado:} considera futuros usuarios, módulos,
		integraciones y cambios en las reglas de negocio.
		
	\end{itemize}
	
	\section*{Procedimiento de selección}
	
	La elección puede realizarse mediante un procedimiento de cinco pasos:
	
	\begin{enumerate}
		
		\item Identificar las funciones principales y los atributos de calidad
		requeridos por la aplicación.
		
		\item Clasificar el tamaño del proyecto, la complejidad del dominio, la
		cantidad de usuarios y el comportamiento de la demanda.
		
		\item Establecer las restricciones de presupuesto, infraestructura,
		tiempo y conocimientos del equipo.
		
		\item Comparar las ventajas y limitaciones de las arquitecturas
		disponibles mediante una matriz de criterios.
		
		\item Crear un prototipo y realizar pruebas de rendimiento, integración y
		mantenimiento antes de aprobar la decisión definitiva.
		
	\end{enumerate}
	
	Este procedimiento evita seleccionar una arquitectura únicamente por
	preferencia tecnológica. También permite documentar las razones de la
	decisión y revisarlas cuando la aplicación cambie.
	
	\section*{Principales arquitecturas}
	
	La arquitectura monolítica integra la interfaz, las reglas de negocio y el
	acceso a datos en una sola aplicación desplegable. Puede ser adecuada para
	proyectos pequeños porque facilita el desarrollo, las pruebas y la
	implementación. IBM indica que el enfoque monolítico continúa siendo una
	alternativa válida para aplicaciones sencillas (Saraswathi, s. f.).
	
	La arquitectura en capas separa responsabilidades como presentación, lógica
	de negocio y acceso a datos. Esta organización es utilizada con frecuencia en
	sistemas empresariales tradicionales porque facilita la distribución de
	responsabilidades. Sin embargo, un cambio puede atravesar varias capas y
	afectar diferentes partes de la aplicación (Microsoft, 2025).
	
	La arquitectura de microservicios divide el sistema en servicios pequeños e
	independientes, organizados alrededor de funciones del negocio. Cada servicio
	puede desarrollarse, desplegarse y escalarse por separado. Esta independencia
	puede beneficiar a sistemas grandes y equipos especializados, pero también
	incrementa la complejidad del monitoreo, la comunicación y la administración
	de datos. Fowler (2015) advierte que los beneficios de los microservicios deben
	compararse con los costos propios de un sistema distribuido.
	
	La arquitectura orientada a eventos permite que los componentes reaccionen
	ante situaciones como una compra, una alerta, una lectura de un sensor o un
	cambio en la información. Resulta útil en aplicaciones de IoT, notificaciones,
	procesamiento en tiempo real y sistemas con comunicación asíncrona.
	
	La arquitectura sin servidor o \textit{serverless} permite ejecutar funciones
	bajo demanda mientras el proveedor de nube administra la infraestructura.
	Puede utilizarse cuando la carga cambia considerablemente o cuando se procesan
	eventos específicos. No obstante, deben evaluarse la dependencia del proveedor,
	los costos acumulados y los límites de ejecución.
	
	\section*{Arquitectura según la aplicación}
	
	La Tabla~\ref{tab:arquitecturas} presenta una relación general entre algunos
	tipos de aplicaciones y las arquitecturas que podrían satisfacer sus
	necesidades. Estas relaciones son recomendaciones y no reglas absolutas.
	
	\begin{table}[H]
		\centering
		\caption{Arquitectura sugerida según el tipo de aplicación}
		\label{tab:arquitecturas}
		\small
		\renewcommand{\arraystretch}{1.20}
		
		\begin{tabularx}{\textwidth}{
				>{\raggedright\arraybackslash}p{3.2cm}
				>{\raggedright\arraybackslash}p{3.4cm}
				>{\raggedright\arraybackslash}X
			}
			\toprule
			\textbf{Tipo de aplicación} &
			\textbf{Arquitectura sugerida} &
			\textbf{Motivo} \\
			\midrule
			
			Pequeña o académica &
			Monolítica o en capas &
			Menor complejidad de desarrollo y despliegue. \\
			
			Empresarial tradicional &
			En capas &
			Separación entre interfaz, reglas de negocio y datos. \\
			
			Grande y modular &
			Microservicios &
			Despliegue y escalamiento independiente de los servicios. \\
			
			IoT o tiempo real &
			Orientada a eventos &
			Procesamiento asíncrono de eventos generados continuamente. \\
			
			Demanda variable &
			Sin servidor &
			Ejecución y escalamiento de funciones bajo demanda. \\
			
			\bottomrule
		\end{tabularx}
		
		\begin{flushleft}
			\footnotesize
			Fuente: elaboración propia con base en Microsoft (2025),
			Saraswathi (s. f.) y Fowler (2015).
		\end{flushleft}
	\end{table}
	
	
	Microsoft (2025) recomienda evaluar si los beneficios de un estilo
	arquitectónico superan sus dificultades dentro del contexto específico de la
	aplicación. Esto confirma que no existe una alternativa universalmente
	superior.
		
	\section*{Observaciones y comentarios}
	
	La arquitectura más compleja no necesariamente es la más adecuada. En
	proyectos pequeños, una solución sencilla y correctamente organizada puede
	ofrecer mejores resultados que una estructura distribuida difícil de
	administrar. También es importante documentar la decisión y revisarla durante
	la evolución del sistema. La arquitectura debe resolver necesidades reales,
	ser comprensible para el equipo y permitir cambios sin aumentar
	innecesariamente los costos.
	
	\section*{Conclusiones}
	
	\begin{enumerate}
		
		\item La arquitectura debe ser seleccionada tomando en cuenta los requisitos, restricciones requeridas por los usuarios.
		
		\item  Una aplicación puede utilizar una arquitectura monolitica o de capas cuando la aplicación es pequeña.
		
		\item No se debe asumir la complejidad de una aplicación, debe siempre realizarse un analisis y tomar en cuenta presupuesto y complejidad.
		
		\item	Elegir una arquitectura de microservicios puede funcionar para aplicaciones más complejas, donde se requiere de escalabilidad e independencia entre sus servicios.
		
		\item	La experiecncia, el presupuesto y el mantenimiento también deben ser evaluados y alineados para poder seleccionar una arquitectura adecuada para una organización.
		
	\end{enumerate}
	
	\section*{Bibliografía}
	
	\begin{itemize}[
		leftmargin=1.27cm,
		label={},
		itemindent=-1.27cm,
		itemsep=6pt
		]
		
		\item Fowler, M. (2015, 1 de julio).
		\textit{Microservice trade-offs}.
		MartinFowler.com.
		\url{https://martinfowler.com/articles/microservice-trade-offs.html}
		
		\item Kazman, R., Echeverria, S., \& Ivers, J. (2023).
		\textit{A holistic view of architecture definition, evolution,
			and analysis} (CMU/SEI-2023-TR-004).
		Software Engineering Institute, Carnegie Mellon University.
		\url{https://doi.org/10.1184/R1/20722402}
		
		\item Microsoft. (2025, 14 de octubre).
		\textit{Architecture styles}.
		Azure Architecture Center.
		\url{https://learn.microsoft.com/en-us/azure/architecture/guide/architecture-styles/}
		
		\item Saraswathi, R. (s. f.).
		\textit{Tipos de arquitecturas de desarrollo de aplicaciones}.
		IBM.
		\url{https://www.ibm.com/es-es/think/topics/application-architecture-types}
	\end{itemize}
	
\end{document}